
\begin{frame}{Case 1: The Disagreement That Followed}

    Northpointe replied that its score was \textbf{equally accurate for both groups}: among
    defendants given the same score, roughly the same fraction re-offended regardless of race.

    \vspace{0.8em}

    \begin{block}{Both sides were reporting true numbers}
        ProPublica measured \textbf{error rates} conditioned on the true outcome
        (who got wrongly flagged). Northpointe measured \textbf{calibration} conditioned on the
        score (what a score means). These are different quantities.
    \end{block}

    \vspace{0.8em}

    \begin{itemize}
        \setlength\itemsep{0.4em}
        \item Chouldechova (2017) and Kleinberg, Mullainathan and Raghavan (2016) showed
              independently that when the two groups have \textbf{different base rates} of the
              outcome, a score cannot satisfy both criteria at once, except in degenerate cases.
        \item So this was not a bug that better engineering would fix. It was a
              \textbf{choice between incompatible definitions of fair} --- and nobody had made it
              explicitly.
    \end{itemize}

    \vspace{0.5em}
    \centering
    \textit{We return to this result in Section 4.}
\end{frame}
